\begin{proof}
\emph{Imported inputs and the finite-order target.}
From~\cite{chen:author2026:time-dependency}*{Theorem~4.1, (4.2), p.~22}, each
deterministic kernel $g_n$ is Stratonovich integrable. Its spectral
hypothesis (1.5) follows from Hypothesis~\ref{H:homog}, as recorded
on p.~3 of the source. The other analytic input is the equal-length
diagram limit in its Lemma~4.4, specified below. We prove the unequal-
length identity from these inputs without invoking a product rule for
multiple Stratonovich integrals.
For the expectation identity needed here it suffices to use the fixed-order
$L^2$ convergence in Theorem~4.1: Cauchy--Schwarz gives $L^1$ convergence
of the product of two regularized factors, since their $L^2$ norms stay
bounded. This is a finite-order statement, not an interchange of infinite
solution series. The cases with a zero order use $X_0=1$.
At positive widths Wick's formula gives zero at odd total order and
the finite pairing sum at even total order; see the source's
(2.20)--(2.21). We justify every term of its noise limit next.

We next justify each individual generalized integral, including unequal
chain lengths. We separate a fixed-noise measure limit from the joint
noise limit imported from the source.

\emph{Green measures at fixed noise widths.}
At fixed positive noise widths, all covariance functions and all finite
Gaussian moments are bounded and continuous in their spatial arguments.
In dimension three the source defines $G_\eta=p_\eta*G$; these positive
measures converge weakly to $G$ and all have mass $s$ at time $s$.
For each fixed strictly ordered time tuple, products of these measures
and their linear pushforwards therefore converge against every bounded
continuous covariance product. The absolute integrals are bounded by
a finite Gaussian-moment constant times the product of the time gaps,
hence by an integrable function on the finite time simplices. Dominated
convergence applies to each finite Wick sum. Applied also to mixed
second moments with two Green widths, it proves the $L^2$ Cauchy
property of every fixed-noise colored integral as $\eta\downarrow0$.
Thus, from now on the fixed-noise integrals use the actual Green
measures. This step takes place before any noise width tends to zero.

\emph{The precise equal-length import.}
From
\cite{chen:author2026:time-dependency}*{Lemma~4.4 and Remark~4.5,
pp.~30--31}, for two chains of the same length $r$, every pairing
$D'\in\Pi_r$ has a finite joint limit when all its vertexwise widths
$(\bar\epsilon,\bar\delta)$ tend to zero. Its regularized covariance
factors are
\[
  \prod_{(j,l)\in D'}
  \gamma^0_{\delta_j\vee\delta_l}(s_j-s_l)
  \gamma_{\epsilon_j+\epsilon_l}(y_j-y_l).
\]
The joint limit permits arbitrary relative rates of all these widths;
it is not a limit with one ratio held fixed. This analytic fact is
imported, not proved here. Its spectral assumption is the source's
(1.5), already checked from Hypothesis~\ref{H:homog}.

\emph{A finite coloring reduces each limit to that import.}
Choose and fix $D\in\Pi_k$ for the original lengths $m,n$.
Assign an independent copy $W^e$ of the source noise to each pair
$e\in D$, and at each vertex use the copy of its pair. For a chain
of length $r\ge1$ write $c(l)$ for this color assignment and
$Y_{r,a}$ for its integral, where $a$ is the collection of its positive
noise widths. At fixed widths it exists by the preceding measure limit.
In the product of two copies of this same length-$r$ chain, define
the doubled color map by
\[
  \widetilde c(j)=c(j),\qquad
  \widetilde c(r+j)=c(j),\qquad 1\le j\le r,
\]
and set
\[
  \Pi_r(c)=\{D'\in\Pi_r:
  \widetilde c(j)=\widetilde c(l)
  \text{ for every }\{j,l\}\in D'\}.
\]
The two copies use widths $a$ and $b$, but the same Gaussian color
fields. Concatenate their vertex widths as
$(\epsilon_j,\delta_j)_{1\le j\le2r}$ and their spacetime variables
as $(\tau_j,\eta_j)_{1\le j\le2r}$. Put
\[
  K_{jl}=\gamma^0_{\delta_j\vee\delta_l}(\tau_j-\tau_l)
  \gamma_{\epsilon_j+\epsilon_l}(\eta_j-\eta_l).
\]
The covariance matrix of the colored evaluations is exactly
$\mathbf1_{\{\widetilde c(j)=\widetilde c(l)\}}K_{jl}$.
Let $H_{r,t,x}(\ud\mathbf s,\ud\mathbf y)$ be the terminal
Green measure on $0<s_1<\cdots<s_r<t$, with kernel product
$\prod_{j=1}^rG(s_{j+1}-s_j,y_{j+1}-y_j)$ and
$(s_{r+1},y_{r+1})=(t,x)$, in the measure convention already stated.
Wick's formula gives exactly
\[
  \begin{split}
    \E[Y_{r,a}Y_{r,b}]&=\sum_{D'\in\Pi_r(c)}B_{D'}(a,b),\\
    B_{D'}(a,b)&=\int H_{r,t,x}(\ud\mathbf s,\ud\mathbf y)
    H_{r,t,x}(\ud\mathbf s',\ud\mathbf y')
    \prod_{\{j,l\}\in D'}K_{jl}.
  \end{split}
\]
An incompatible pairing has a zero color indicator; for a compatible
pairing every indicator is one. The remaining integrand has no colors:
it is precisely the one-field fixed-diagram integral with these $2r$
vertex widths. The one-field import already uses the same covariance
function on all its $r$ edges, so repetition adds no new hypothesis.
For example, if $r=2$ and both chain vertices have one color, the
four-vertex covariance product in the expectation is
\[
  K_{12}K_{34}+K_{13}K_{24}+K_{14}K_{23}.
\]
These are exactly the three uncolored diagrams in $\Pi_2$, including
the self-pairing on each copy. None is discarded.
Reversing each terminal chain as specified below gives its initial-chain
diagram of Lemma~4.4, merely permuting the vertex widths. Each
$B_{D'}$ therefore has a finite limit $b_{D'}$, independent of their
relative rates. Define the finite sum
$Q_{r,c}=\sum_{D'\in\Pi_r(c)}b_{D'}$. Consequently,
\[
  \lim_{a,b\to0}\E[Y_{r,a}Y_{r,b}]=Q_{r,c},\qquad
  \lim_{a\to0}\E Y_{r,a}^2=\lim_{b\to0}\E Y_{r,b}^2=Q_{r,c}.
\]
It follows that
\[
  \lim_{a,b\to0}\|Y_{r,a}-Y_{r,b}\|_2^2
  =Q_{r,c}+Q_{r,c}-2Q_{r,c}=0.
\]
Completeness of $L^2$ gives a unique limit $Y_r$. Set $Y_0=1$
for an empty chain.

\emph{The original unequal chains and their reversal.}
For the original two chains the product converges in $L^1$, since
\[
  \begin{split}
    \|Y_{m,a}Y_{n,b}-Y_mY_n\|_1
    &\le\|Y_{m,a}-Y_m\|_2\|Y_{n,b}\|_2\\
    &\quad+\|Y_m\|_2\|Y_{n,b}-Y_n\|_2\longrightarrow0.
  \end{split}
\]
In that product each color occurs exactly twice in total, even when
both occurrences lie on one chain. Thus, exactly one pairing of equal
colors is possible: the prescribed $D$. Its regularized expectation
is the terminal-chain integral for $D$, and has the finite limit
$\E[Y_mY_n]$. The reversal below carries it to $J_{m,n,D^{\rm rev}}$.
Since $D\mapsto D^{\rm rev}$ is a bijection of pairings, every
integral in~\eqref{E:mixed-diagram} is obtained this way.
The regularized integrands are nonnegative, since the Green measures,
temporal covariances and smoothed spatial covariances are nonnegative.
Hence, every limiting $J_{m,n,D}$ is nonnegative as well. This proves
finiteness and nonnegativity of each diagram, without interchanging
an infinite series or asserting a pointwise value of a singular kernel.

More precisely, define the involution on vertex labels by
\[
  \rho(j)=m+1-j\quad(1\le j\le m),\qquad
  \rho(m+j)=m+n+1-j\quad(1\le j\le n),
\]
and set $D^{\rm rev}=\{\{\rho(a),\rho(b)\}:\{a,b\}\in D\}$.
Since $\rho^2$ is the identity, this is a bijection of $\Pi_k$.

More explicitly, if $(s_j,y_j)$ are the first chain's terminal
variables, set
\[
  s'_j=t-s_{m+1-j},\qquad y'_j=x-y_{m+1-j},
  \qquad 1\le j\le m,
\]
and use the analogous length-$n$ transformation on the second chain.
The new variables remain strictly increasing in time and have the
initial root $(s'_0,y'_0)=(0,0)$ on each chain. For example, the
terminal factor $G(t-s_m,x-y_m)$ becomes $G(s'_1,y'_1)$;
every other Green factor becomes its corresponding successive
increment. This also specifies the equal-length reversal used above.
Both time and spatial covariance differences change sign; the covariance
kernels are even. The resulting permutation of the $2k$ labels is a
bijection of $\Pi_k$. At positive widths the uncolored product is
the finite sum of these individual terminal diagrams. Its $L^1$
limit and the individual colored limits therefore give
\eqref{E:mixed-moment}, including its
nonnegativity. These are exactly the two chains in the integrand of
the source's (3.6), with $n_1=m$, $n_2=n$, restricted to the finite
horizons. The same transformations apply to the source's regularized
integrals before taking their defining limits.

\emph{The separate solution-series passage.}
Proposition~\ref{P:mild-construction} below identifies the series
of the $X_n$ as a Stratonovich mild solution. Its fixed-$t$ summability
follows independently from Lemma~\ref{L:crude-corrected} and the
vertexwise version in Lemma~\ref{L:vertex-crude}. Their proofs use
the source's fixed-order inputs, not the present mixed expansion.
Thus, this use of the later construction is not circular.
In particular, $U_N=\sum_{n=0}^N X_n$ converges to $u$ in $L^2$.
Consequently,
\[
  |\E U_N^2-\E u^2|
  \le\|U_N-u\|_2(\|U_N\|_2+\|u\|_2)\longrightarrow0.
\]
The finite expansion of $\E U_N^2$ is a sum of the nonnegative mixed
expectations already computed. Its increasing limit is therefore the
full double sum, independently of its enumeration. Regrouping by total
order gives~\eqref{E:full-second}. This proves Proposition~\ref{P:mixed}.
\end{proof}
